% =====================================================================
%  THÈME 3 — Proportions stœchiométriques
%  Niveau DIFFICILE — Corrigé
% =====================================================================
\documentclass[11pt,a4paper]{article}
\usepackage[utf8]{inputenc}
\usepackage[T1]{fontenc}
\usepackage{lmodern}
\usepackage[french]{babel}
\usepackage{amsmath,amssymb,mathtools}
\usepackage{icomma}
\usepackage[version=4]{mhchem}
\usepackage{siunitx}
\sisetup{output-decimal-marker={,},per-mode=symbol}
\usepackage{xcolor}
\usepackage{tikz}
\usepackage[most]{tcolorbox}
\usepackage{enumitem}
\usepackage{array}
\usepackage{fancyhdr}
\usepackage[a4paper,margin=2.1cm,headheight=26pt,headsep=8pt]{geometry}

\definecolor{primary}{RGB}{142,93,168}
\definecolor{primarydark}{RGB}{82,45,108}
\definecolor{corrcol}{RGB}{176,110,20}
\newcommand{\nq}[1]{n_{\text{#1}}}
\newcommand{\Vm}{V_{\mathrm m}}

\pagestyle{fancy}\fancyhf{}
\fancyhead[L]{\small\textcolor{primarydark}{\textbf{Fiche 5} — Thème 3 : Proportions}}
\fancyhead[R]{\small\textcolor{corrcol}{\textsc{Corrigé — Difficile}}}
\fancyfoot[C]{\small\thepage}
\renewcommand{\headrulewidth}{0.8pt}
\renewcommand{\headrule}{\hbox to\headwidth{\color{primary}\leaders\hrule height \headrulewidth\hfill}}

\newtcolorbox[auto counter]{cor}[1][]{enhanced,breakable,boxrule=1pt,arc=3pt,
  left=3mm,right=3mm,top=2.4mm,bottom=2.4mm,colback=corrcol!6,colframe=corrcol,
  fonttitle=\bfseries,coltitle=white,
  title={Corrigé~\thetcbcounter\ \ifx\relax#1\relax\relax\else\textmd{--- \itshape #1}\fi}}

\newcommand{\rep}[1]{\par\smallskip\noindent\textbf{\color{corrcol}$\Rightarrow$ #1}}

\setlength{\parindent}{0pt}
\setlength{\parskip}{3pt}

\begin{document}

\begin{center}
{\LARGE\bfseries\color{primarydark}Thème 3 — Proportions}\\[1pt]
{\large\color{corrcol}\textbf{Corrigé — Niveau Difficile}}\\
{\color{primary}\rule{0.6\linewidth}{1pt}}
\end{center}

\begin{cor}[bilan massique complet d'une synthèse]
Rapport de référence : $\dfrac{\nq{Al2(SO4)3}}{1}=0{,}100$.\\
\textbf{a)} $\nq{Al(OH)3}=2\times0{,}100=0{,}200$~mol $\Rightarrow m=0{,}200\times78=\boxed{\SI{15,6}{\gram}}$.\\
\phantom{a)} $\nq{H2SO4}=3\times0{,}100=0{,}300$~mol $\Rightarrow m=0{,}300\times98=\boxed{\SI{29,4}{\gram}}$.\\
\textbf{b)} $\nq{H2O}=6\times0{,}100=0{,}600$~mol $\Rightarrow m=0{,}600\times18=\boxed{\SI{10,8}{\gram}}$.\\
\textbf{c)} $m(\ce{Al2(SO4)3})=0{,}100\times342=\SI{34,2}{\gram}$.
\[
m_{\text{réactifs}}=15{,}6+29{,}4=\SI{45,0}{\gram};\qquad m_{\text{produits}}=34{,}2+10{,}8=\SI{45,0}{\gram}. \quad\checkmark
\]
\rep{La masse se conserve (Lavoisier) : c'est un excellent contrôle de cohérence.}
\end{cor}

\begin{cor}[remonter d'un volume de gaz à une masse de réactif]
\textbf{a)} $\nq{C2H2}=\dfrac{V}{\Vm}=\dfrac{60}{24}=\SI{2,5}{\mole}$ (ici $\Vm=\SI{24}{\litre\per\mole}$).\\
\textbf{b)} Ratio \ce{CaC2}:\ce{C2H2}$=1:1$ : $\nq{CaC2}=\SI{2,5}{\mole}$.\\
\textbf{c)} $m(\ce{CaC2})=n\times M = 2{,}5\times64=\boxed{\SI{160}{\gram}}$.
\rep{Attention : ici $\Vm=24$ (à \SI{298}{\kelvin}), pas $22{,}4$ (qui vaut aux TPN).}
\end{cor}

\begin{cor}[un airbag : de la masse au volume de gaz]
\textbf{a)} $\nq{NaN3}=\dfrac{13{,}0}{65}=\SI{0,200}{\mole}$.\\
\textbf{b)} Ratio \ce{NaN3}:\ce{N2}$=2:3$ : $\nq{N2}=\nq{NaN3}\times\dfrac{3}{2}=0{,}200\times1{,}5=\SI{0,300}{\mole}$.\\
\textbf{c)} $V(\ce{N2})=n\Vm=0{,}300\times22{,}4=\boxed{\SI{6,72}{\litre}}$.
\rep{Un airbag réel utilise plusieurs dizaines de grammes de \ce{NaN3} pour \(\sim\)\SI{60}{\litre} de gaz.}
\end{cor}

\begin{cor}[combustion du propane : masse, volume et Lavoisier]
\textbf{a)} $\nq{C3H8}=\dfrac{44}{44}=\SI{1,0}{\mole}$.\\
\textbf{b)} $\nq{O2}=5\times1{,}0=5{,}0$~mol $\Rightarrow V=5{,}0\times22{,}4=\boxed{\SI{112}{\litre}}$.\\
\textbf{c)} $\nq{CO2}=3\times1{,}0=3{,}0$~mol $\Rightarrow m=3{,}0\times44=\boxed{\SI{132}{\gram}}$ ;\quad
$\nq{H2O}=4\times1{,}0=4{,}0$~mol $\Rightarrow m=4{,}0\times18=\boxed{\SI{72}{\gram}}$.\\
\textbf{d)} $m_{\text{réactifs}}=\underbrace{44}_{\ce{C3H8}}+\underbrace{5{,}0\times32}_{\ce{O2}=160}=\SI{204}{\gram}$ ;\quad
$m_{\text{produits}}=132+72=\SI{204}{\gram}$. \checkmark
\rep{Masse conservée : le bilan est cohérent.}
\end{cor}

\end{document}
